\section{Khảo sát các giải pháp phần mềm hiện có}
\label{sec:solutions}

\subsection{Các hệ thống PACS/VNA thương mại}

Thị trường phần mềm quản lý và lưu trữ ảnh y khoa từ lâu do một số nhà cung cấp
lớn chi phối. GE Healthcare, Siemens Healthineers và Philips đều phát triển các
hệ thống PACS gắn liền với dòng thiết bị chẩn đoán hình ảnh của chính mình.
Fujifilm với dòng sản phẩm Synapse, cùng Sectra và Agfa HealthCare, cung cấp các
nền tảng PACS và VNA hướng tới quy mô bệnh viện và mạng lưới nhiều cơ sở.

Điểm chung của nhóm này là mức độ hoàn chỉnh cao. Hệ thống được cung cấp trọn
gói, kèm dịch vụ triển khai, đào tạo và bảo trì, tích hợp sẵn với HIS và RIS
theo các chuẩn HL7 và DICOM, đồng thời phối hợp chặt chẽ với máy chụp cùng hãng.
Đổi lại, chi phí đầu tư và phí duy trì hằng năm ở mức cao, mã nguồn đóng và cấu
trúc dữ liệu nội bộ thường không được công bố, khiến cơ sở y tế rơi vào tình
trạng phụ thuộc nhà cung cấp cả về nâng cấp, mở rộng lẫn di chuyển dữ liệu về
sau.

\subsection{Các hệ thống PACS mã nguồn mở}

\textbf{Orthanc} là một PACS mã nguồn mở nhẹ với giấy phép thoáng. Hệ thống chạy
từ một tệp nhị phân duy nhất, không bắt buộc phải có cơ sở dữ liệu bên ngoài, hỗ
trợ đồng thời giao thức DICOM cổ điển và DICOMweb, đồng thời tích hợp bộ thông
dịch Lua cho phép định nghĩa các hành vi tự động khi có ảnh mới đến. Nhờ yêu cầu
tài nguyên thấp, Orthanc phù hợp cả khi đóng vai trò PACS đầy đủ tại một cơ sở
nhỏ lẫn khi làm cổng tích hợp đặt cạnh một PACS sẵn có tại cơ sở lớn
\cite{jodogne2018orthanc}.

Điểm đáng chú ý của Orthanc là hệ sinh thái plugin, cho phép bổ sung đúng những
năng lực mà lõi còn thiếu mà không phải sửa lõi \cite{orthancbook}. Plugin
PostgreSQL thay SQLite mặc định bằng một cơ sở dữ liệu chịu tải cao làm chỉ mục, và
cho phép nhiều phiên bản Orthanc dùng chung chỉ mục. Plugin lưu trữ đối tượng ghi
dữ liệu điểm ảnh lên kho tương thích S3, Azure Blob hoặc Google Cloud Storage thay
vì hệ thống tệp cục bộ. Plugin DICOMweb cung cấp WADO-RS, STOW-RS và QIDO-RS. Plugin
phân quyền nâng cao chuyển mọi quyết định cấp quyền cho một dịch vụ web bên ngoài,
và dự án orthanc-auth-service hiện thực sẵn dịch vụ này trên nền máy chủ định danh
Keycloak, kèm cơ chế chia sẻ ca chụp bằng token \cite{orthancauth,keycloak}. Tổ hợp
này giải quyết được một phần đáng kể các điểm yếu về quản trị nêu ở
Mục~\ref{sec:solutions-analysis}, nhưng vẫn là việc lắp ghép cần được thiết kế và
kiểm thử.

\textbf{dcm4che} là bộ công cụ DICOM viết bằng Java, đã trưởng thành qua nhiều
năm phát triển và được xem là một trong những hiện thực DICOM mã nguồn mở đầy đủ
nhất. Trên nền bộ công cụ này, \textbf{dcm4chee-arc} cung cấp một kho lưu trữ ảnh
ở quy mô doanh nghiệp, hỗ trợ phong phú các dịch vụ DICOM và DICOMweb cùng các
cơ chế quản trị dữ liệu. Bù lại, hệ thống đòi hỏi hạ tầng nặng hơn đáng kể, gồm
máy chủ ứng dụng Java và cơ sở dữ liệu riêng, cũng như nhân sự quen thuộc với hệ
sinh thái Java.

\subsection{Các trình xem ảnh y khoa}

\textbf{OHIF Viewer}, do Open Health Imaging Foundation duy trì \cite{ziegler2020ohif}, là trình xem mã
nguồn mở chạy hoàn toàn trên trình duyệt. Đặc tính không cần cài đặt tại máy
trạm (zero-footprint) giúp loại bỏ khâu triển khai và cập nhật phần mềm trên
từng máy tính, một lợi thế rõ rệt tại các cơ sở có năng lực quản trị công nghệ
thông tin hạn chế. Kiến trúc mở rộng theo mô-đun cho phép bổ sung công cụ đo
đạc, bố cục hiển thị hay lớp phủ thông tin ngay trong giao diện đọc ảnh tiêu
chuẩn. Từ phiên bản 3, OHIF dựa trên Cornerstone3D, hỗ trợ dựng hình bằng GPU và tải
ảnh lũy tiến (Mục~\ref{sec:transfersyntax}).

\textbf{Cornerstone.js} là thư viện hiển thị ảnh y khoa trên nền web, đảm nhận
việc giải mã, dựng hình và tương tác với dữ liệu DICOM. Đây chính là lớp nền mà
OHIF Viewer xây dựng bên trên, và cũng có thể được dùng trực tiếp khi cần một
giao diện xem ảnh gọn nhẹ, gắn chặt với ứng dụng sẵn có.

\textbf{Weasis} là trình xem chạy trên máy trạm, viết bằng Java, hỗ trợ nhiều
phương thức chụp và có thể khởi chạy trực tiếp từ liên kết do PACS sinh ra. Ưu
điểm của trình xem desktop là khai thác tốt tài nguyên máy cục bộ, nhưng đi kèm
với đó là yêu cầu cài đặt và duy trì phần mềm trên từng máy.

\subsection{Các thành phần mã nguồn mở chuyên biệt}

Ngoài PACS và trình xem, một số dự án mã nguồn mở giải quyết đúng từng khâu mà
đề tài cần, và là tham chiếu trực tiếp cho thiết kế ở Mục~\ref{sec:architecture}.

\textbf{KHEOPS} là nền tảng chia sẻ ảnh DICOM, tổ chức ca chụp theo ``album'' và
cấp quyền truy cập bằng token gắn với một album hoặc một ca chụp, có thể giới hạn
quyền và thời hạn; dữ liệu được truy cập qua DICOMweb và mở bằng OHIF hoặc Weasis
\cite{kheops}. Mô hình token phạm vi hẹp này khớp với kịch bản gửi một ca chụp cho
bác sĩ tuyến trên mà không cấp cho họ quyền vào toàn bộ kho ảnh.

\textbf{Karnak}, của cùng nhóm phát triển, là cổng DICOM chuyên khử định danh và
định tuyến ảnh theo các hồ sơ của DICOM PS3.15, với giao diện cấu hình quy tắc khử
định danh cho từng dự án \cite{karnak}. Đây là tham chiếu cho khâu khử định danh
tại ranh giới của bộ điều phối.

\textbf{MONAI Deploy Informatics Gateway} tiếp nhận dữ liệu qua DICOM DIMSE,
DICOMweb và FHIR, gom các đối tượng thành ca chụp, phát sự kiện kích hoạt quy trình
xử lý và gửi kết quả ngược về PACS \cite{monaiig}. Về chức năng, nó gần như trùng
với bộ điều phối trong mô hình thành phần, và là ứng viên nên được đánh giá trước
khi tự xây dựng, đặc biệt khi các thành phần phân tích ảnh tự động được bổ sung ở
các giai đoạn sau. Đổi lại, nó cần thêm hàng đợi thông điệp và kho đối tượng, nặng
so với thiết bị trọn gói ở bối cảnh C.

\subsection{Các nền tảng ảnh y khoa trên đám mây}
\label{sec:cloud}

Các nhà cung cấp đám mây lớn đều có dịch vụ quản lý ảnh y khoa: AWS HealthImaging
\cite{awshealthimaging}, kho DICOM của Google Cloud Healthcare API
\cite{gcphealthcare}, và dịch vụ DICOM của Azure Health Data Services, vốn phát
triển từ dự án mã nguồn mở Medical Imaging Server for DICOM của Microsoft
\cite{msdicom}. Cả ba đều cung cấp DICOMweb, tự mở rộng quy mô và tính phí theo
dung lượng và lưu lượng. Chúng thay thế tầng lưu trữ, không thay thế trình xem
hay quy trình đọc ảnh.

Bài viết của AWS về xây dựng PACS gốc đám mây \cite{awspacs2026} đáng được phân
tích kỹ vì nó đi tới cùng kết luận với đề tài từ hướng ngược lại. Kiến trúc được đề
xuất là mô hình lai trục--nan: mỗi bệnh viện giữ một máy chủ VNA cục bộ tiếp nhận
ảnh qua C-STORE và phục vụ bác sĩ tại chỗ ở tốc độ mạng nội bộ, còn ảnh được sao
chép bất đồng bộ lên kho đối tượng trung tâm qua kênh riêng mà không chặn luồng lâm
sàng. Một cơ sở dữ liệu trung tâm chỉ chứa siêu dữ liệu, gồm chỉ mục bệnh nhân
chung, vị trí của từng ca chụp và trạng thái đồng bộ, cho phép bác sĩ ở bệnh viện
này tìm và mở ca chụp thực hiện ở bệnh viện khác; trình xem tự định tuyến yêu cầu
tới bản cục bộ nếu còn trong bộ đệm, hoặc tới kho trung tâm nếu không. Bài viết
dựa trên quan sát rằng tần suất truy cập ảnh cao trong vài tháng đầu, giảm mạnh sau
6--12 tháng và hiếm khi được truy cập sau 2--3 năm để đề xuất phân tầng lưu trữ theo
vòng đời, từ lớp truy cập thường xuyên tới lớp lưu trữ sâu cho dữ liệu trên năm năm,
và khuyến nghị mô hình lai thay vì thuần đám mây khi kết nối mạng hạn chế hoặc chỉ
phụ thuộc một nhà mạng.

Đối chiếu với đề tài: (i) nguyên tắc ghi cục bộ trước, đồng bộ sau
(Mục~\ref{sec:localfirst}) trùng khớp, cho thấy đây không chỉ là lựa chọn bắt buộc
của cơ sở thiếu tài nguyên mà còn là thực hành được khuyến nghị ở quy mô lớn; (ii)
tách chỉ mục khỏi điểm ảnh và phân tầng theo vòng đời là hai điểm đề tài tiếp thu,
trình bày ở Mục~\ref{sec:tiering} và Mục~\ref{sec:hub}; (iii) khác biệt nằm ở giả
định. Bài viết giả định kênh kết nối riêng ổn định, đội ngũ vận hành, và khả năng đặt
dữ liệu tại một vùng đám mây trong lãnh thổ quốc gia; ở Việt Nam, dữ liệu đưa ra
ngoài lãnh thổ phải qua đánh giá tác động chuyển dữ liệu xuyên biên giới
(Mục~\ref{sec:legal}), và ở bối cảnh C kênh kết nối chính là điểm yếu nhất. Vì vậy
đề tài đặt trung tâm ở hạ tầng trong nước và chỉ gửi ra ngoài bệnh viện bản sao lưu
đã mã hoá và bản đã khử định danh.


\subsection{Bối cảnh trong nước}

Tại Việt Nam, chủ trương chuyển đổi số ngành y tế và lộ trình triển khai bệnh án
điện tử đang thúc đẩy các bệnh viện thay thế phim in bằng hệ thống lưu trữ và
truyền tải hình ảnh. Thông tư 13/2025/TT-BYT đặt thời hạn triển khai hồ sơ bệnh án
điện tử cho các bệnh viện chậm nhất ngày 30/9/2025 và cho các cơ sở khám bệnh, chữa
bệnh khác chậm nhất ngày 31/12/2026 \cite{tt13_2025}; vì kết quả cận lâm sàng, trong
đó có chẩn đoán hình ảnh, là một phần của hồ sơ bệnh án, lộ trình này tạo áp lực
trực tiếp lên việc số hoá và lưu trữ ảnh ở chính nhóm cơ sở chưa có PACS. Nhiều cơ sở lớn lựa chọn giải pháp thương mại nhập khẩu
hoặc dịch vụ do doanh nghiệp trong nước cung cấp, trong khi tuyến huyện và các
phòng khám phần lớn vẫn gặp trở ngại về chi phí đầu tư và nhân lực vận hành. Đây
là khoảng trống mà một giải pháp xây dựng trên nền tảng mã nguồn mở có thể lấp
đầy.

\section{Phân tích các giải pháp hiện có}
\label{sec:solutions-analysis}

\subsection{Điểm mạnh}

Các hệ thống PACS thương mại mang lại sự hoàn chỉnh và tính ổn định đã được kiểm
chứng qua nhiều năm vận hành lâm sàng. Người mua nhận được một sản phẩm trọn gói
cùng cam kết hỗ trợ kỹ thuật, quy trình tích hợp với HIS và RIS đã được chuẩn
hoá, và trách nhiệm khắc phục sự cố thuộc về một đầu mối duy nhất. Với bệnh viện
có lưu lượng chụp lớn và ngân sách tương xứng, đây là phương án giảm thiểu rủi ro
triển khai.

Nhóm mã nguồn mở lại mạnh ở những khía cạnh khác. Chi phí cấp phép gần như bằng
không cho phép dồn nguồn lực vào hạ tầng và con người thay vì bản quyền phần mềm.
Mã nguồn công khai giúp kiểm tra được cách dữ liệu được lưu trữ và xử lý, đồng
thời cho phép tuỳ biến sâu theo quy trình thực tế của từng khoa chẩn đoán hình
ảnh. Việc tuân thủ chặt các chuẩn DICOM và DICOMweb giữ cho dữ liệu luôn ở định
dạng mở, nhờ đó có thể chuyển đổi giữa các hệ thống mà không bị ràng buộc. Riêng
Orthanc và OHIF còn nổi bật ở yêu cầu tài nguyên thấp, đủ để vận hành trên một
máy chủ phổ thông tại cơ sở tuyến dưới.

\subsection{Điểm yếu và khoảng trống}

\textbf{Chi phí và cấp phép.} Các hệ thống thương mại được định giá cho những
bệnh viện dồi dào ngân sách, theo hình thức mua bản quyền kèm phí bảo trì hằng
năm hoặc tính theo dung lượng lưu trữ. Với bệnh viện tuyến huyện và phòng khám
tư nhân quy mô nhỏ, mức chi phí này vượt quá khả năng chi trả, trong khi đây lại
chính là nhóm cơ sở cần một hệ thống quản lý ảnh có tổ chức nhất.

\textbf{Phụ thuộc nhà cung cấp.} Khi hệ thống đóng kín và cấu trúc lưu trữ không
được công bố, mọi thay đổi về sau đều phải đi qua nhà cung cấp ban đầu. Việc mở
rộng thêm chức năng, kết nối với hệ thống của bên thứ ba hay di chuyển toàn bộ
kho ảnh sang nền tảng khác đều trở nên tốn kém, và cơ sở y tế mất phần lớn quyền
chủ động đối với hạ tầng của chính mình.

\textbf{Khả năng tuỳ biến.} Quy trình làm việc tại mỗi khoa chẩn đoán hình ảnh
có những khác biệt nhất định, từ cách phân luồng ca chụp đến biểu mẫu kết quả.
Sản phẩm thương mại chỉ cho phép cấu hình trong khuôn khổ nhà sản xuất định sẵn,
còn những yêu cầu nằm ngoài khuôn khổ đó thường phải chờ bản cập nhật hoặc trả
phí phát triển riêng.

\textbf{Hạ tầng và nhân lực vận hành.} Ở chiều ngược lại, các giải pháp mã nguồn
mở quy mô doanh nghiệp như dcm4chee-arc đòi hỏi đội ngũ có kinh nghiệm quản trị
máy chủ, cơ sở dữ liệu và mạng. Nhiều cơ sở tuyến dưới không có sẵn nhân lực như
vậy, nên một hệ thống tuy miễn phí nhưng phức tạp vẫn có thể trở thành gánh nặng
vận hành.

\textbf{Triển khai tại cơ sở hạn chế tài nguyên.} Phần lớn giải pháp thương mại
mặc định rằng cơ sở đã có máy chủ đủ mạnh, đường truyền ổn định và bộ phận công
nghệ thông tin thường trực. Những nơi không đáp ứng được các điều kiện đó gần
như không có phương án phù hợp, và khả năng vận hành ở chế độ hạn chế khi mạng
gián đoạn hiếm khi được tính đến.

\textbf{Chủ quyền dữ liệu.} Các dịch vụ lưu trữ ảnh trên đám mây đòi hỏi dữ liệu
hình ảnh rời khỏi phạm vi cơ sở y tế, làm phát sinh những ràng buộc mà
Mục~\ref{sec:legal} cho thấy không đơn thuần mang tính hợp đồng theo pháp luật
Việt Nam. Khả năng lưu trữ và xử lý hoàn toàn tại chỗ vì vậy là một yêu cầu
thực chất chứ không phải tuỳ chọn.

\textbf{Gánh nặng tích hợp của các thành phần mã nguồn mở.} Hệ sinh thái mã
nguồn mở tránh được chi phí cấp phép nhưng chỉ cung cấp các thành phần rời rạc,
không phải một giải pháp hoàn chỉnh. Orthanc, OHIF và các thư viện liên quan,
xét riêng từng thành phần, đều đã trưởng thành; nhưng ghép chúng thành một hệ
thống có xác thực người dùng, phân quyền theo vai trò, nhật ký kiểm toán, giám
sát và hành vi được định nghĩa rõ ràng khi xảy ra sự cố là một khối lượng công
việc kỹ thuật đáng kể. Chính phần lắp ghép này tạo nên đóng góp của đề tài.

\begin{table}[H]
\centering
\small
\begin{tabularx}{\textwidth}{@{}X X X@{}}
\toprule
\textbf{Giải pháp} & \textbf{Điểm mạnh} & \textbf{Điểm yếu} \\
\midrule
PACS thương mại (GE, Siemens, Philips, Fujifilm, Sectra, Agfa) & Hoàn chỉnh; hỗ trợ chính hãng; tích hợp sâu với thiết bị và với HIS/RIS & Chi phí đầu tư và duy trì cao; mã nguồn đóng; phụ thuộc nhà cung cấp; khó tuỳ biến \\
Orthanc & Nhẹ, một tệp nhị phân; hỗ trợ cả DICOM cổ điển lẫn DICOMweb; scripting Lua; giấy phép thoáng & Thiếu sẵn các chức năng quản trị cấp bệnh viện; phải tự bổ sung xác thực, phân quyền và giám sát \\
dcm4che/dcm4chee-arc & Trưởng thành; đầy đủ dịch vụ DICOM; phù hợp quy mô doanh nghiệp & Hạ tầng nặng; cấu hình phức tạp; đòi hỏi nhân lực quản trị Java và cơ sở dữ liệu \\
OHIF Viewer & Chạy trên trình duyệt, không cần cài đặt tại máy trạm; kiến trúc mở rộng theo mô-đun & Phụ thuộc năng lực máy khách và đường truyền; cần tuỳ biến để khớp quy trình đọc ảnh cụ thể \\
Weasis & Trình xem desktop giàu chức năng; khai thác tốt tài nguyên máy cục bộ & Phải cài đặt và cập nhật trên từng máy trạm \\
Cornerstone.js & Thư viện hiển thị linh hoạt, nhúng được vào ứng dụng web sẵn có & Chỉ là thư viện, không phải trình xem hoàn chỉnh; phải tự xây dựng giao diện \\
Plugin Orthanc (PostgreSQL, lưu trữ đối tượng, phân quyền) & Bổ sung chỉ mục chịu tải, lưu trữ S3, phân quyền qua Keycloak mà không sửa lõi & Mỗi plugin là một điểm cấu hình và nâng cấp riêng; cần kiểm thử tổ hợp \\
KHEOPS & Chia sẻ ca chụp bằng token phạm vi hẹp; DICOMweb; tích hợp OHIF và Weasis & Thêm một hệ thống cần vận hành; không tiếp nhận trực tiếp từ máy chụp \\
Karnak & Khử định danh theo hồ sơ PS3.15, cấu hình qua giao diện & Chỉ giải quyết khâu khử định danh và định tuyến \\
MONAI Deploy Informatics Gateway & Tiếp nhận DIMSE, DICOMweb, FHIR; phát sự kiện; trả kết quả về PACS & Hướng tới quy trình AI; phụ thuộc hàng đợi thông điệp và kho đối tượng \\
Dịch vụ đám mây (AWS HealthImaging, Google Cloud Healthcare API, Azure DICOM) & Không phải vận hành hạ tầng lưu trữ; tự mở rộng; DICOMweb & Dữ liệu rời khỏi bệnh viện; ràng buộc chuyển dữ liệu xuyên biên giới; phụ thuộc đường truyền và nhà cung cấp \\
\bottomrule
\end{tabularx}
\caption{Phân tích so sánh các giải pháp quản lý và hiển thị ảnh y khoa hiện có.}
\label{tab:solutions}
\end{table}
